\chapter{光的波动性：衍射与干涉}\label{ch:g12:light-as-wave}

将肥皂膜朝向窗户倾斜：膜上的颜色不断滑动，
而肥皂和水本身并无颜色。用激光照射一根头发：
身后的墙上出现的不是细如发丝的影子，而是一道明亮的光带，
其中布满暗纹，宽度是头发的数千倍。
沿直线传播的光线（\cref{ch:g10:refraction}）无法解释
这两种视觉现象。本章将光提升为完整的波——波峰到波峰
相距约半微米——继而将这一思想付诸应用：\omterm{def:g11:color-light-sources:lamps}{激光器}
与裸露的墙面可用来测量头发丝的粗细、一个\omterm{def:g12:mechanical-waves:wavelength}{波长}以及 CD 碟片上的刻纹。

\section{光是一种波}

\begin{definition}[光波]\label{def:g12:light-as-wave:lightwave}
\omterm{def:g10:light-spectra:white}{单色光}是一种纯色，棱镜无法将其进一步分开
（\cref{ch:g10:light-spectra}），是一种\omterm{def:g10:signals-and-waves:period}{周期性的}\emph{光
波}\index{光波}：它在真空中传播的速率为 $c = \qty{3.00e8}{m/s}$，
并由其\omterm{def:g10:signals-and-waves:frequency}{频率} $f$ 或其\emph{真空中的
波长}\index{波长!光的} $\lambda = c/f$，即
相邻波峰之间的距离，见\cref{ch:g12:mechanical-waves}。人眼
能\omterm{def:g11:color-light-sources:perception}{感知}约 \qty{400}{nm}（紫）至 \qty{800}{nm}（红）之间的光；
\omterm{def:g10:light-spectra:white}{白光}则混合了整个范围内的光。
\end{definition}

\begin{example}[数量级]\label{ex:g12:light-as-wave:orders}
氦氖\omterm{def:g11:color-light-sources:lamps}{激光器}发出的红光，其\omterm{def:g12:mechanical-waves:wavelength}{波长}为 $\lambda = \qty{633}{nm} =
\qty{6.33e-7}{m}$，故 $f = c/\lambda \approx \qty{4.7e14}{Hz}$：
任何电子元件都无法如此快速地计数——本章中的每个\omterm{def:g12:mechanical-waves:wavelength}{波长}都将
借助几何方法测量。作为参照，一根头发（约 \qty{70}{\micro m}）
的宽度相当于一百个\omterm{def:g12:mechanical-waves:wavelength}{波长}。什么在波动？不是空气：是\omterm{def:g11:electric-gravitational-fields:efield}{电场}
与\omterm{def:g11:magnetic-fields:bfield}{磁场}，其严格描述见本科二年级分册；下文只
关心\omterm{def:g12:mechanical-waves:wavelength}{波长}。
\end{example}

\section{衍射}

\begin{definition}[衍射]\label{def:g12:light-as-wave:diffraction}
当波遇到尺寸为 $a$ 且与其\omterm{def:g12:mechanical-waves:wavelength}{波长}相当的孔径或障碍物时，
它便会扩展进入射线理论中本应为暗的区域。
这种扩展就是\emph{衍射}\index{衍射}——
当 $\lambda/a$ 很小时可以忽略，$a$ 减小至接近 $\lambda$ 时则占主导。
\end{definition}

\begin{omfigure}
\begin{tikzpicture}[font=\small]
  \foreach \x in {-2.2,-1.75,...,-0.4}
    \draw[omDef, thick] (\x,-1.6) -- (\x,1.6);
  \draw[very thick] (0,-1.6) -- (0,-0.3) (0,0.3) -- (0,1.6);
  \foreach \r in {0.45,0.9,...,2.5}
    \draw[omDef, thick] (-75:\r) arc (-75:75:\r);
  \draw[Stealth-Stealth] (0.14,-0.3) -- (0.14,0.3)
    node[pos=0.28, right] {$a$};
\end{tikzpicture}

{\small 平面\omterm{def:g12:mechanical-waves:wavefront}{波前}（每个波峰一条，彼此相距 $\lambda$）射到一条
宽度为 $a$、与 $\lambda$ 相当的狭缝上：波散开，即\omterm{def:g12:light-as-wave:diffraction}{衍射}。}
\end{omfigure}

\begin{proposition}[扩散的角半宽]\label{prop:g12:light-as-wave:halfwidth}
宽度为 $a$ 的狭缝，以\omterm{def:g12:mechanical-waves:wavelength}{波长} $\lambda < a$ 的光照射时，
光会围绕入射方向扩散，其角半宽
为 $\theta$（弧度）；距离为 $D$ 的屏幕上呈现中央亮带，
亮带宽度 $L$ 为 $D \tan\theta \approx D\,\theta$ 的两倍（小
角度近似）：
\[
\theta \approx \frac{\lambda}{a}, \qquad L = \frac{2\lambda D}{a}.
\]
\end{proposition}
\admitted

\begin{remark}[反过来看公式]\label{rem:g12:light-as-wave:honest}
狭缝越\emph{窄}，扩散越\emph{宽}，
与射线预言恰好相反。严格推导（对狭缝上各子波求和）
见本科二年级分册；此处公式仅是
一种测量工具。
\end{remark}

\begin{example}[激光通过狭缝]\label{ex:g12:light-as-wave:slit}
红色\omterm{def:g11:color-light-sources:lamps}{激光器}，$\lambda = \qty{633}{nm}$；狭缝 $a = \qty{0.10}{mm}$；
屏幕位于 $D = \qty{2.00}{m}$ 处。于是 $\theta = \num{6.33e-7} /
\num{1.0e-4} = \qty{6.3e-3}{rad}$——约合三分之一度，而
墙上的中央亮带宽度为 $L = 2 \times \num{6.33e-7} \times 2.00 /
\num{1.0e-4} \approx \qty{2.5}{cm}$，两侧为暗线及
较弱的旁侧亮带。狭缝减半，亮带则加倍。
\end{example}

\begin{omfigure}
\begin{tikzpicture}[font=\small]
  \fill[gray!55] (-3.2,-0.22) rectangle (-2.2,0.22);
  \node[above] at (-2.7,0.25) {激光器};
  \draw[omThm, very thick] (-2.2,0) -- (0,0);
  \draw[very thick] (0,-1.5) -- (0,-0.1) (0,0.1) -- (0,1.5);
  \node[below] at (0,-1.5) {狭缝 $a$};
  \draw[dashed] (0,0) -- (4.6,0);
  \draw[omThm] (0,0) -- (4.6,0.9) (0,0) -- (4.6,-0.9);
  \draw (1.5,0) arc (0:11:1.5) node[right, xshift=1pt] {$\theta$};
  \draw[very thick] (4.6,-1.85) -- (4.6,1.85);
  \draw[Stealth-Stealth] (4.9,-0.9) -- (4.9,0.9)
    node[midway, right] {$L$};
  \draw[Stealth-Stealth] (0,-2.1) -- (4.6,-2.1)
    node[midway, below] {$D$};
\end{tikzpicture}

{\small 单缝实验：半角 $\theta
\approx \lambda/a$ 的光锥在墙上照出中央亮带，其宽度为 $L = 2\lambda D/a$。}
\end{omfigure}

\begin{proposition}[障碍物的衍射如同狭缝]\label{prop:g12:light-as-wave:hair}
宽度为 $d$ 的细不透明障碍物——如金属丝、头发——
在入射光束之外，与宽度为 $d$ 的狭缝产生相同的图样，
（巴比涅原理，推导见本科二年级分册）。
\end{proposition}
\admitted

\begin{method}[测量微小宽度]\label{met:g12:light-as-wave:hairwidth}
\begin{enumerate}
  \item 将头发横跨在激光束上，
        并在测得距离 $D$（几米）处放置屏幕。
  \item 测量中央亮带的宽度 $L$，即两侧暗纹之间的距离。
  \item 由 $L = 2\lambda D/d$ 可得 $d = 2\lambda D/L$——头发越细，
        图样越大。
\end{enumerate}
\end{method}

\begin{remark}[衍射限制每一种仪器]\label{rem:g12:light-as-wave:limits}
光通过宽度为 $a$ 的孔径进入照相机、望远镜或眼睛，
因此即使是理想\omterm{def:g11:lenses-and-eye:lens}{透镜}，也会把一个点模糊成角尺寸约为 $\lambda/a$ 的光斑：
较近的细节会融合在一起。可见光显微镜因此几乎无法分辨
$\lambda \approx
\qty{0.5}{\micro m}$ 以下的结构（细菌可分辨，病毒不能）；望远镜反射镜
也因同样原因而做得很宽。
\end{remark}

\section{干涉}

\begin{definition}[相干波源，光程差]\label{def:g12:light-as-wave:coherent}
两个波源在\omterm{def:g10:signals-and-waves:frequency}{频率}相同、振动步调一致时，称为\emph{相干波源}\index{相干波源}——
实际上，它们是同一列波经分束得到的两份副本。对于
两者都能到达的点 $M$，
其\emph{光程差}%
\index{光程差}为 $\delta = S_2M - S_1M$。
\end{definition}

\begin{proposition}[干涉条件]\label{prop:g12:light-as-wave:conditions}
两列\omterm{def:g12:mechanical-waves:wavelength}{波长}为 $\lambda$ 的相干波在重叠处相加——
即它们发生\emph{干涉}\index{干涉}：
\begin{itemize}
  \item $\delta = k\lambda$（$k$ 为整数）：步调一致，相互加强——
        \emph{相长干涉}\index{相长干涉}（亮纹）；
  \item $\delta = (k + \tfrac12)\lambda$：波峰与波谷相遇，\omterm{def:g10:inertia:compensate}{相互抵消}
        ——\emph{相消干涉}\index{相消干涉}（暗纹）。
\end{itemize}
\end{proposition}

\begin{proof}
将\omterm{def:g10:signals-and-waves:period}{周期}性波平移整数个\omterm{def:g12:mechanical-waves:wavelength}{波长}，结果不变：波峰
总与波峰叠加。再移动半个\omterm{def:g12:mechanical-waves:wavelength}{波长}，则波峰对准波谷：
两列等幅波\omterm{def:g10:inertia:compensate}{相互抵消}。
\end{proof}

\begin{example}[杨氏双缝实验]\label{ex:g12:light-as-wave:young}
在不透明薄板上穿出两条细缝 $S_1$、$S_2$，相距不到一毫米，
并用同一台\omterm{def:g11:color-light-sources:lamps}{激光器}照亮两缝，得到两个\omterm{def:g12:light-as-wave:coherent}{相干
波源}。在相距几米的屏上，重叠区域呈条纹状：
在中心处 $\delta = 0$，为亮纹；沿屏向上，$\delta$
逐渐增大，经过 $\lambda/2$（暗纹）、$\lambda$（亮纹）\,\dots\ 后，
出现暗纹，即光与光相加反而抵消之处。
\end{example}

\begin{omfigure}
\begin{tikzpicture}[font=\small]
  \draw[very thick] (0,-1.7) -- (0,-0.72) (0,-0.48) -- (0,0.48)
    (0,0.72) -- (0,1.7);
  \fill (0,0.6) circle (1.5pt) node[left=2pt] {$S_1$} (0,-0.6) circle (1.5pt) node[left=2pt] {$S_2$};
  \draw[Stealth-Stealth] (-0.75,-0.6) -- (-0.75,0.6)
    node[midway, left] {$a$};
  \draw[very thick] (5.0,-2.0) -- (5.0,2.0);
  \draw[omDef] (0,0.6) -- (5.0,1.3) (0,-0.6) -- (5.0,1.3);
  \fill (5.0,1.3) circle (1.5pt) node[above right] {$M$};
  \draw[dashed] (0,0) -- (5.0,0) node[below right] {$O$};
  \draw[Stealth-Stealth] (5.55,0.04) -- (5.55,1.26) node[midway, right] {$x$};
  \draw[dashed] ([shift={(5.0,1.3)}]188:5.05) arc (188:200.8:5.05);
  \draw[omProp, very thick] (0,-0.6) -- (0.281,-0.493)
    node[below right=-2pt] {$\delta$};
  \foreach \y in {-1.8,-1.35,...,1.8}
    \fill[omExo] (5.06,\y-0.1) rectangle (5.2,\y+0.1);
\end{tikzpicture}

{\small 杨氏双缝：到 $M$ 的\omterm{def:g12:light-as-wave:coherent}{光程差}为 $\delta = S_2M -
S_1M$；亮条纹（灰色刻线）位于 $\delta = k\lambda$ 处。}
\end{omfigure}

\begin{definition}[条纹与条纹间距]\label{def:g12:light-as-wave:fringes}
图样中明暗交替的条带是它的
\emph{干涉条纹}\index{干涉条纹}；相邻亮纹中心之间的距离
$i$ 称为\emph{条纹间距}%
\index{条纹间距}。
\end{definition}

\begin{proposition}[条纹间距]\label{prop:g12:light-as-wave:interfringe}
若两缝相距 $a$，屏幕位于距离 $D \gg a$ 处，
屏上距离中心为 $x$ 的点（中心为 $O$），其
\omterm{def:g12:light-as-wave:coherent}{光程差} $\delta \approx a\,x/D$（几何关系在本阶段不加证明地承认）；由
\cref{prop:g12:light-as-wave:conditions}可知，亮纹位于
$x_k = k\,\lambda D/a$ 处，等间距，\omterm{def:g12:light-as-wave:fringes}{条纹间距}为
\[
i = \frac{\lambda D}{a}.
\]
已知 $\lambda = \qty{633}{nm}$，$a = \qty{0.20}{mm}$，$D =
\qty{2.0}{m}$，则 $i = \qty{6.3}{mm}$——该装置把 $\lambda$ 放大了
$D/a = 10^4$ 倍。其几何推导在
本科二年级分册中给出。
\end{proposition}
\admitted

\begin{omfigure}
\begin{tikzpicture}
\begin{axis}[omaxis, width=8.6cm, height=4.4cm,
    xlabel={$x$（\unit{mm}）}, ylabel={光强},
    xmin=-11, xmax=11, ymin=0, ymax=1.3, ytick=\empty]
  \addplot[omDef, thick, domain=-10.5:10.5, samples=211]
    {cos(deg(pi*x/3))^2 * exp(-x*x/50)};
  \draw[Stealth-Stealth] (axis cs:0,1.1) -- (axis cs:3,1.1)
    node[midway, above] {$i$};
\end{axis}
\end{tikzpicture}

{\small 杨氏实验中屏幕上的光强分布，\omterm{def:g12:light-as-wave:fringes}{干涉条纹}间距为
$i = \lambda D/a$，外包络线仅为示意。}
\end{omfigure}

\begin{method}[测量波长]\label{met:g12:light-as-wave:lambda}
\begin{enumerate}
  \item 让未知激光通过间距已知的双缝，缝距为
        $a$；将屏幕置于测得的距离 $D$ 处。
  \item 测量 $10$ 个\omterm{def:g12:light-as-wave:fringes}{条纹间距}的总跨度，再除以 $10$，得到\omterm{def:g12:light-as-wave:fringes}{条纹间距}
        $i$，尺子误差因此缩小为原来的十分之一。
  \item 于是 $\lambda = i\,a/D$。
\end{enumerate}
\end{method}

\begin{remark}[两盏台灯为何不产生干涉]\label{rem:g12:light-as-wave:lamps}
两盏台灯照亮同一面墙，不会出现\omterm{def:g12:light-as-wave:fringes}{干涉条纹}：每盏灯发出的是
短促、彼此无关的波列，其相对相位每秒钟随机跳变十亿次，
条纹也随之飞快移动，
眼睛便将它们平均为均匀的光。\omterm{prop:g12:light-as-wave:conditions}{干涉}需要相干性，
即一束波被分成两束，如杨氏双缝所做的那样；
\omterm{def:g11:color-light-sources:lamps}{激光器}则是现代的捷径，它发出的是干净而长的单一波列。
\end{remark}

\section{干涉描绘的颜色}

\begin{definition}[虹彩]\label{def:g12:light-as-wave:iridescence}
随观察角度而变化的颜色——肥皂泡、油膜、
孔雀羽毛、CD 的背面——称为\emph{虹彩}%
\index{虹彩}：并非源于色素，而是\omterm{def:g10:light-spectra:white}{白光}的\omterm{prop:g12:light-as-wave:conditions}{干涉}，
使某些\omterm{def:g12:mechanical-waves:wavelength}{波长}得到加强，另一些则被抵消。
\end{definition}

\begin{example}[肥皂膜与油膜]\label{ex:g12:light-as-wave:soap}
光照射到肥皂膜上时，一部分从前表面反射，
另一部分从后表面反射：两束相干光的\omterm{def:g12:light-as-wave:coherent}{光程差}
由膜的厚度和观察角度决定。每种厚度都会使某些\omterm{def:g12:mechanical-waves:wavelength}{波长}得到加强，
另一些则被抵消：薄膜反射的颜色随液体流下、膜变薄而不断变化。
湿润沥青路面上的油膜，
也是同样的道理。
\end{example}

\begin{remark}[单色光与白光]\label{rem:g12:light-as-wave:white}
在\omterm{def:g10:light-spectra:white}{白光}中，每种\omterm{def:g12:mechanical-waves:wavelength}{波长}都会形成各自的\omterm{def:g12:light-as-wave:fringes}{干涉条纹}，
其间距 $i = \lambda D/a$ 各不相同：中央处各色光都明亮，形成一条白色条纹；
离中央几条条纹之外，各色图样开始分离，边缘呈现\omterm{def:g12:light-as-wave:iridescence}{虹彩}，
随后变得模糊。\omterm{def:g10:light-spectra:white}{单色光}则给出清晰、无穷尽的\omterm{def:g12:light-as-wave:fringes}{干涉条纹}，
这也是精密测量中都用\omterm{def:g11:color-light-sources:lamps}{激光器}的原因。
\end{remark}

\section{练习}

\begin{exercise}[$\star$]\label{exo:g12:light-as-wave:1}
计算红光（$\lambda = \qty{700}{nm}$）与
紫光（\qty{400}{nm}）的\omterm{def:g10:signals-and-waves:frequency}{频率}。哪一种\omterm{def:g10:signals-and-waves:frequency}{频率}更高？
\end{exercise}

\begin{exercise}[$\star$]\label{exo:g12:light-as-wave:2}
一束绿色激光（$\lambda = \qty{532}{nm}$）穿过宽度为
$a = \qty{0.20}{mm}$ 的单缝：求 $\theta$，
再求离屏幕 \qty{1.5}{m} 处中央亮带的宽度。
\end{exercise}

\begin{exercise}[$\star$]\label{exo:g12:light-as-wave:3}
门的宽度为 \qty{0.80}{m}：求 $\lambda/a$，分别针对
\omterm{def:g10:signals-and-waves:frequency}{频率}为 \qty{425}{Hz} 的\omterm{def:g10:signals-and-waves:sound}{声音}（$v = \qty{340}{m/s}$）与 \qty{550}{nm} 的光。
为什么能听到拐角后的\omterm{def:g10:signals-and-waves:sound}{声音}，却看不到拐角后的光？
\end{exercise}

\begin{exercise}[$\star$]\label{exo:g12:light-as-wave:4}
由同一发生器驱动的两个扬声器\omterm{def:g12:mechanical-waves:phase}{同相}发出
\omterm{def:g12:mechanical-waves:wavelength}{波长}为 \qty{2.0}{cm} 的\omterm{def:g10:signals-and-waves:sound}{声音}。在 $M$ 点的\omterm{def:g12:light-as-wave:coherent}{光程差}为 \qty{6.0}{cm}；在
$N$ 点为 \qty{5.0}{cm}。两处各为响还是静？请说明理由。
\end{exercise}

\begin{exercise}[$\star$]\label{exo:g12:light-as-wave:5}
杨氏双缝：$\lambda = \qty{633}{nm}$，$a = \qty{0.25}{mm}$，
$D = \qty{1.8}{m}$。求\omterm{def:g12:light-as-wave:fringes}{条纹间距}，屏幕中央
\qty{2.0}{cm} 内可容纳多少条亮纹？
\end{exercise}

\begin{exercise}[$\star\star$]\label{exo:g12:light-as-wave:6}
在 $\lambda = \qty{650}{nm}$ 的光束中放一根头发，屏幕距离为 \qty{1.6}{m}，
测得中央亮带宽度为 \qty{2.8}{cm}。求头发的直径。
\end{exercise}

\begin{exercise}[$\star\star$]\label{exo:g12:light-as-wave:7}
未知激光通过杨氏双缝（$a = \qty{0.20}{mm}$，
$D = \qty{1.50}{m}$）：十个\omterm{def:g12:light-as-wave:fringes}{条纹间距}的总跨度为 \qty{4.7}{cm}。求
$\lambda$ 及其颜色。为什么要测量十个而不是一个？
\end{exercise}

\begin{exercise}[$\star\star$]\label{exo:g12:light-as-wave:8}
两盏相同的台灯照亮同一面墙，光束相互重叠，
却始终不出现\omterm{def:g12:light-as-wave:fringes}{干涉条纹}。哪一个相干条件
（\cref{def:g12:light-as-wave:coherent}）不满足？杨氏装置
又是如何弥补这一点的？
\end{exercise}

\begin{exercise}[$\star\star$]\label{exo:g12:light-as-wave:9}
在双缝图样中（$a = \qty{0.30}{mm}$，$D = \qty{1.4}{m}$），
第一条与第七条亮纹中心相距 \qty{18.0}{mm}：
求\omterm{def:g12:light-as-wave:fringes}{条纹间距}，再求 $\lambda$。
\end{exercise}

\begin{exercise}[$\star\star$]\label{exo:g12:light-as-wave:10}
CD 的磁道构成间距为 $d = \qty{1.6}{\micro m}$ 的光栅。
取 $d \sin\theta = k\lambda$（$k$ 为整数），求
$\lambda = \qty{633}{nm}$ 时的光束方向。最高级次 $k$ 为多少？
\end{exercise}

\begin{exercise}[$\star\star$]\label{exo:g12:light-as-wave:11}
口径为 $a = \qty{10}{cm}$ 的望远镜，\omterm{def:g12:mechanical-waves:wavelength}{波长} $\lambda =
\qty{550}{nm}$ 时，求其\omterm{def:g12:light-as-wave:diffraction}{衍射}模糊角 $\theta \approx
\lambda/a$。能否分辨相距 \qty{1.5}{m} 的两盏车头灯，
它们距望远镜 \qty{100}{km}？
\end{exercise}

\begin{exercise}[$\star\star\star$]\label{exo:g12:light-as-wave:12}
竖直放置的肥皂膜在流动变薄，下部较厚，并呈现
水平彩色条纹向下缓慢移动。试解释为何有颜色，
为何出现条纹，为何会移动。（不需计算。）
\end{exercise}

\begin{exercise}[$\star\star\star$]\label{exo:g12:light-as-wave:13}
\omterm{def:g10:light-spectra:white}{白光}下的杨氏双缝，$a = \qty{0.20}{mm}$，$D = \qty{1.5}{m}$：
求紫光（\qty{400}{nm}）与红光
（\qty{750}{nm}）的\omterm{def:g12:light-as-wave:fringes}{条纹间距}，再推断屏幕上
中央、离中央几毫米处及远处分别呈现什么。
\end{exercise}

\begin{exercise}[$\star\star\star$]\label{exo:g12:light-as-wave:14}
一颗侦察\omterm{def:g10:universal-gravitation:satellite}{卫星}携带口径为 \qty{2.4}{m} 的镜子，\omterm{def:g10:universal-gravitation:satellite}{轨道}高度为
\qty{250}{km}。取 $\theta \approx \lambda/a$，\omterm{def:g12:mechanical-waves:wavelength}{波长}为 \qty{550}{nm}，
求其能分辨的最小地面细节，并评判其
“能隔着你的肩膀读报”这一说法。
\end{exercise}

\begin{exercise}[$\star\star\star$]\label{exo:g12:light-as-wave:15}
杨氏双缝装置在空气中显示 $i = \qty{4.0}{mm}$，随后整体浸入
水中（$n = 1.33$）。回顾\cref{ch:g10:refraction}，\omterm{def:g10:signals-and-waves:frequency}{频率}是否
改变？\omterm{def:g12:mechanical-waves:wavelength}{波长}呢？求新的\omterm{def:g12:light-as-wave:fringes}{条纹间距}。
\end{exercise}

\section{问题：用光测量}

\begin{problem}[{周末问题——用光测量：激光器、直尺与
空白墙面变成微米测量工坊，用以测定狭缝、一根头发丝、
未知波长以及两张银色圆盘的刻槽}]
\label{pb:g12:light-as-wave:1}
你的器材：红色\omterm{def:g11:color-light-sources:lamps}{激光器}（$\lambda = \qty{633}{nm}$，取自标签），
一台\omterm{def:g12:mechanical-waves:wavelength}{波长}未知的绿色\omterm{def:g11:color-light-sources:lamps}{激光器}，一条已标定的单缝
$a = \qty{100}{\micro m}$，一条间距为 $a' =
\qty{0.250}{mm}$ 的双缝、一把卷尺、两根头发、一张 CD 和一张 DVD。

\medskip
\noindent\textbf{第一部分——用已知单缝标定。}
红色\omterm{def:g11:color-light-sources:lamps}{激光器}、单缝与墙面，墙面距离为 $D = \qty{3.00}{m}$。
\begin{enumerate}
  \item 求\omterm{def:g12:light-as-wave:diffraction}{衍射}光束的角半宽 $\theta$。
  \item 证明中央亮带宽度为 $L = 2\lambda D/a$；
        并求其值。
  \item 若测得 $L = \qty{3.7}{cm}$，求其相对
        偏差，该方法是否得到验证？
  \item 第二条单缝宽度减半，为 \qty{50}{\micro m}，预测其
        $L$。陈述缝宽与图样宽度之间的关系规律。
  \item 说明第 2 问中取 $\tan\theta \approx \theta$ 这一步的理由：比较
        该角度下 $\theta$ 与 $\tan\theta$ 的大小。
\end{enumerate}

\medskip
\noindent\textbf{第二部分——头发丝。}\omterm{def:g11:color-light-sources:lamps}{激光器}与墙面不变，
用绷紧的头发丝代替单缝。
\begin{enumerate}[resume]
  \item 哪一条在本阶段不加证明地承认的结论，使你能将单缝公式用于
        头发丝？该结论的内容是什么？
  \item 若中央亮带宽 $L = \qty{5.4}{cm}$，求
        头发丝的直径 $d$。
  \item 你的直尺测得 $L$ 的误差为 $\pm\qty{0.2}{cm}$，
        由此给出 $d$ 的范围，写成 $d \pm \Delta d$。
  \item 你朋友的头发给出 $L = \qty{7.6}{cm}$，其直径为多少？
        谁的头发更细？为何头发越细，中央亮带反而越宽？
  \item 能否用台灯代替\omterm{def:g11:color-light-sources:lamps}{激光器}？给出两条理由。
\end{enumerate}

\medskip
\noindent\textbf{第三部分——绿色\omterm{def:g11:color-light-sources:lamps}{激光器}的\omterm{def:g12:mechanical-waves:wavelength}{波长}。}绿色
\omterm{def:g11:color-light-sources:lamps}{激光器}、双缝，墙面距离为 $D = \qty{2.00}{m}$。
\begin{enumerate}[resume]
  \item 为何两条缝必须由\emph{同一}台\omterm{def:g11:color-light-sources:lamps}{激光器}照亮？
  \item 为何无论\omterm{def:g12:mechanical-waves:wavelength}{波长}如何，中央条纹都是亮的？
  \item 十个\omterm{def:g12:light-as-wave:fringes}{条纹间距}共宽 \qty{4.26}{cm}。求 $i$，并说明
        为何要一次测量十个。
  \item 由此推出绿色\omterm{def:g11:color-light-sources:lamps}{激光器}的\omterm{def:g12:mechanical-waves:wavelength}{波长}。
  \item 厂家标签标注为 $\qty{532 \pm 10}{nm}$，二者是否一致？
        你预计会看到什么颜色？
\end{enumerate}

\medskip
\noindent\textbf{第四部分——读取银色圆盘。}红色\omterm{def:g11:color-light-sources:lamps}{激光器}
垂直入射到 CD 上；其磁道起到反射光栅的作用，
遵循（在本阶段不加证明地承认的）$d \sin\theta = k\lambda$。墙面位于 $D =
\qty{1.00}{m}$ 处，每束一级光束形成一个光斑，
位于距中央光斑 $x$ 处。
\begin{enumerate}[resume]
  \item 对于 CD，$x = \qty{43}{cm}$，求 $\theta$。为何
        第一部分的小角近似现在为何不再允许使用？
  \item 推出 CD 的磁道间距 $d$。
  \item DVD 的一级光斑移至 $x = \qty{1.65}{m}$ 处：
        求其 $\theta$ 与磁道间距。
  \item 比较两种间距。为何 DVD 比 CD 存储的信息更多？
  \item 合上笔记本：列出本周末所测得的长度，
        并用一句话说明：是什么让光成为一把通用的尺子。
\end{enumerate}
\end{problem}
